\cotefile{83}
\scanpage{178}
l'homomorphisme voulu de
\[
\deleted{\unsure{f}}\quad E_1\mathbin{*}_{\lambda}E_2\quad\text{dans}\quad E_1^\circ * E_2^\circ.
\]
Par passage à la limite \unsure{etc.}, on en conclut l'injectivité si les $F_i$ sont plats.

Il faut définir un $\lambda$-\unsure{semi-hom.} \inserted{(\uline{canonique}) sur} $\mathop{*}E_i$, donc il faut définir
\[
\mathop{*}^{i}E_i\longrightarrow\bigotimes E_i^\circ
\]
qui sur $\bigotimes E_i^\circ$ prenne la valeur $\lambda\,\mathrm{id}$.
Il faut donc trouver un $\lambda$-accouplement \unsure{canon. $i$-linéaire} dans $\bigotimes E_i^\circ$.

On songe \deleted{\unsure{tout}} de suite à :
\[
f=\bigotimes u_i:\ \bigotimes E_i\longrightarrow\bigotimes E_i^\circ
\]
\[
f(x_1,\ldots x_n)=u_1(x_1)\otimes\cdots\otimes u_n(x_n)
\]
\[
\begin{aligned}
f(x_1+\alpha_1,x_2\ldots x_n)
 &=\bigl(u_1(x_1)+\lambda\alpha_1\bigr)\otimes u_2(x_2)\cdots\otimes u_n(x_n)\\
 &=f(x_1,\ldots x_n)+\underbrace{\lambda\alpha_1}_{\deleted{\unsure{\alpha_1}}}\otimes u_2(x_2)\cdots u_n(x_n)
\end{aligned}
\]
\margin{mais il faudrait remplacer $\bigotimes E_i$ par son \unsure{im.} --- ou bien ne pas \unsure{considérer} dans les données
\diagram{\begin{tikzcd}[ampersand replacement=\&]F_i\arrow[r,"\unsure{v_i}"]\&E_i\end{tikzcd}}
\unsure{pour celles qui sont} injectives mais seulement \unsure{se donner un} $u_i:E_i\unsure{\to|-}F_i$ \unsure{avec} $u_i\unsure{v_i}=\lambda\,\mathrm{id}_{F_i}$.}

prendre donc comme application associée
\[
f^\circ(\alpha_1,\ldots,\alpha_n)=\lambda\alpha_1\otimes\cdots\otimes\alpha_n
\]
de sorte que les termes \unsure{en nb premiers} \inserted{\unsure{d'ordre}}\margin{177}
